\begin{figure}[tbp]
\centering
\begin{tikzpicture}[
  >=Latex,
  party/.style={circle,draw,fill=white,minimum size=6.5mm,inner sep=0pt,font=\small},
  push/.style={->,very thick},
  every node/.style={font=\small}
]
  \begin{scope}[xshift=-3.8cm]
    \coordinate (p1) at (0,2.05);
    \coordinate (p2) at (1.78,1.03);
    \coordinate (p3) at (1.78,-1.03);
    \coordinate (p4) at (0,-2.05);
    \coordinate (p5) at (-1.78,-1.03);
    \coordinate (p6) at (-1.78,1.03);

    % The three four-party supports used by the displayed loop.
    \fill[black!7] (p1) -- (p2) -- (p3) -- (p4) -- cycle;
    \fill[pattern=north east lines,pattern color=black!35]
      (p2) -- (p3) -- (p4) -- (p5) -- cycle;
    \fill[pattern=dots,pattern color=black!38]
      (p1) -- (p4) -- (p5) -- (p6) -- cycle;

    \foreach \n in {1,...,6} \node[party] at (p\n) {\n};
    \draw[push] (p1) -- node[right,xshift=1mm] {$M_A(1,3)$} (p3);
    \draw[push] (p3) -- node[above,yshift=1mm] {$M_B(3,5)$} (p5);
    \draw[push] (p5) -- node[left,xshift=-1mm] {$M_C(5,1)$} (p1);
    \node[align=left,anchor=north west,font=\scriptsize] at (-2.55,-2.38) {
      \tikz\fill[black!12] (0,0) rectangle (0.34,0.16);\; $A=\{1,2,3,4\}$\\[-1pt]
      \tikz\fill[pattern=north east lines,pattern color=black!45] (0,0) rectangle (0.34,0.16);\; $B=\{2,3,4,5\}$\\[-1pt]
      \tikz\fill[pattern=dots,pattern color=black!45] (0,0) rectangle (0.34,0.16);\; $C=\{1,4,5,6\}$
    };
  \end{scope}

  \draw[->,thick] (-0.55,0) -- (0.25,0)
    node[midway,above,font=\scriptsize] {compose};

  \begin{scope}[xshift=2.65cm]
    \draw[->] (-1.45,0) -- (1.55,0) node[right] {$P_1$};
    \draw[->] (0,-1.4) -- (0,1.45);
    \fill (0.42,0.52) circle (1.8pt) node[above right] {$v$};
    \fill (-0.72,0.78) circle (1.8pt) node[above left] {$Hv$};
    \draw[->,thick,bend left=18] (0.46,0.55) to
      node[below,font=\small] {$H$} (-0.68,0.82);
    \node[draw,rounded corners,align=center,inner sep=4pt] at (0,-2.0)
      {$F_1H=HF_1$\\[-1pt]
       \scriptsize compatible gauges centralize every loop};
  \end{scope}
\end{tikzpicture}
\caption{Operator pushing on minimum supports, shown for six parties.
Each \((m+1)\)-party support identifies the local Weyl-label planes of its
parties through the bijective projections of the support theorem; overlapping
supports compose to a loop whose holonomy \(H=M_CM_BM_A\) is a linear map of
the base plane, carrying a label \(v\) of the base plane to \(Hv\).
An atlas gauge is compatible precisely when it intertwines every transition;
for one state over a prime field, its base block centralizes every loop
holonomy.}
\label{fig:operator-holonomy}
\end{figure}
